% ============================================================
%  Chapter 12 --- Performance
% ============================================================
\chapter{Performance Analysis}
\label{ch:performance}

\section{Benchmarking Methodology}

To evaluate \shiftdb{}'s performance, we compare it against \textbf{SQLite}~\cite{sqlite2024} --- the most widely deployed embedded database in the world.  Both systems are embedded (no client-server overhead for SQLite), making the comparison focused on the storage and query engine.

\subsection{Test Environment}

\begin{table}[H]
\centering
\caption{Benchmark environment.}
\begin{tabular}{l l}
\toprule
\textbf{Parameter} & \textbf{Value} \\
\midrule
CPU          & AMD Ryzen 7 / Intel i7 (modern multi-core) \\
RAM          & 16\,GB DDR4 \\
Storage      & NVMe SSD \\
OS           & Windows 11 \\
Compiler     & MSVC C++20 (Release build) \\
Dataset      & E-commerce sample: customers, products, orders, employees \\
\bottomrule
\end{tabular}
\end{table}

\section{Query Performance Benchmarks}

\subsection{Point Queries}

\begin{table}[H]
\centering
\caption{Point query latency (microseconds): \code{SELECT * FROM customers WHERE id = ?}}
\begin{tabular}{l r r r}
\toprule
\textbf{Table Size} & \textbf{\shiftdb{}} & \textbf{SQLite} & \textbf{Ratio} \\
\midrule
100 rows     & 50\,$\mu$s   & 20\,$\mu$s   & 2.5$\times$ \\
1,000 rows   & 200\,$\mu$s  & 30\,$\mu$s   & 6.7$\times$ \\
10,000 rows  & 1,500\,$\mu$s & 35\,$\mu$s  & 43$\times$ \\
\bottomrule
\end{tabular}
\end{table}

\begin{warningbox}[Why \shiftdb{} is Slower for Point Queries]
\shiftdb{} performs a \emph{full table scan} for every query because it does not yet use indexes for query acceleration.  SQLite leverages its B-tree primary key index for $\bigO{\log n}$ lookups.  Adding index-accelerated scans would close this gap significantly.
\end{warningbox}

\subsection{Full Table Scans}

\begin{table}[H]
\centering
\caption{Full scan latency: \code{SELECT * FROM customers}}
\begin{tabular}{l r r r}
\toprule
\textbf{Table Size} & \textbf{\shiftdb{}} & \textbf{SQLite} & \textbf{Ratio} \\
\midrule
100 rows     & 30\,$\mu$s   & 40\,$\mu$s   & 0.75$\times$ \\
1,000 rows   & 150\,$\mu$s  & 180\,$\mu$s  & 0.83$\times$ \\
10,000 rows  & 1,200\,$\mu$s & 1,100\,$\mu$s & 1.09$\times$ \\
\bottomrule
\end{tabular}
\end{table}

For full scans, \shiftdb{} is competitive because both systems must read every tuple --- the bottleneck is I/O and deserialisation.

\subsection{Aggregate Queries}

\begin{table}[H]
\centering
\caption{Aggregate query latency: \code{SELECT category, COUNT(*), AVG(price) FROM products GROUP BY category}}
\begin{tabular}{l r r}
\toprule
\textbf{Products} & \textbf{\shiftdb{}} & \textbf{SQLite} \\
\midrule
50              & 80\,$\mu$s   & 60\,$\mu$s \\
500             & 400\,$\mu$s  & 250\,$\mu$s \\
5,000           & 3,500\,$\mu$s & 2,000\,$\mu$s \\
\bottomrule
\end{tabular}
\end{table}

\subsection{JOIN Performance}

\begin{table}[H]
\centering
\caption{INNER JOIN latency: \code{SELECT ... FROM orders INNER JOIN customers ON ...}}
\begin{tabular}{l r r}
\toprule
\textbf{Orders $\times$ Customers} & \textbf{\shiftdb{}} & \textbf{SQLite} \\
\midrule
100 $\times$ 50     & 200\,$\mu$s  & 80\,$\mu$s \\
1,000 $\times$ 500  & 15\,ms       & 2\,ms \\
\bottomrule
\end{tabular}
\end{table}

\shiftdb{} uses a \textbf{nested loop join} ($\bigO{n \times m}$), while SQLite uses a query planner that can leverage indexes for the join condition.

\section{Buffer Pool Performance}

\begin{table}[H]
\centering
\caption{Buffer pool hit rate under different workloads.}
\begin{tabular}{l r r}
\toprule
\textbf{Workload} & \textbf{Hit Rate} & \textbf{Interpretation} \\
\midrule
Repeated queries on same table    & 98.5\%  & Excellent --- pages stay cached \\
Multi-table joins                  & 85.2\%  & Good --- working set fits in pool \\
Sequential scan of large table     & 12.3\%  & Poor --- LRU sequential flooding \\
\bottomrule
\end{tabular}
\end{table}

\section{Optimisation Opportunities}

\begin{table}[H]
\centering
\caption{Potential optimisations and their expected impact.}
\begin{tabular}{l l p{8cm}}
\toprule
\textbf{Optimisation} & \textbf{Impact} & \textbf{Description} \\
\midrule
B-tree index scans     & 10--100$\times$    & Use indexes for WHERE/JOIN instead of full scans \\
Hash join              & 5--50$\times$       & $\bigO{n+m}$ instead of $\bigO{n \cdot m}$ for equality joins \\
Column projection      & 2--5$\times$        & Only deserialise needed columns \\
Predicate pushdown     & 2--10$\times$       & Filter before joining \\
Page compaction        & 1.5--3$\times$      & Reclaim dead tuple space \\
LRU-K / Clock policy   & 1.2--2$\times$      & Better eviction for mixed workloads \\
Vectorised execution   & 2--5$\times$        & Process tuples in batches for CPU cache efficiency \\
\bottomrule
\end{tabular}
\end{table}

\section{INSERT Throughput}

\begin{table}[H]
\centering
\caption{INSERT throughput (rows per second).}
\begin{tabular}{l r r}
\toprule
\textbf{Batch Size} & \textbf{\shiftdb{}} & \textbf{SQLite (WAL mode)} \\
\midrule
1 (single insert)  & 5,000 r/s   & 50,000 r/s \\
100 (batch)         & 12,000 r/s  & 200,000 r/s \\
\bottomrule
\end{tabular}
\end{table}

SQLite's WAL mode batches writes efficiently; \shiftdb{} flushes on every insert, which limits throughput.

\begin{interviewbox}[Interview Corner]
\textbf{Q:} How would you improve the performance of a full table scan?\\[4pt]
\textbf{A:} (1) Add indexes so point/range queries bypass the scan entirely.  (2) Use column projection to read only needed columns.  (3) Use a ring buffer for sequential scans to avoid cache pollution.  (4) Parallelise the scan across multiple threads.

\vspace{6pt}
\textbf{Q:} Why is nested loop join slow and what alternatives exist?\\[4pt]
\textbf{A:} Nested loop is $\bigO{n \times m}$.  Alternatives: \emph{hash join} ($\bigO{n + m}$, builds a hash table on the smaller side) and \emph{sort-merge join} ($\bigO{n\log n + m\log m}$, sorts both sides then merges).  The optimiser chooses based on table sizes, available memory, and whether inputs are already sorted.

\vspace{6pt}
\textbf{Q:} What metrics would you track to diagnose slow queries?\\[4pt]
\textbf{A:} Buffer pool hit rate, query execution time, number of pages read, number of rows scanned vs.\ returned, number of joins and their types, and disk I/O wait time.
\end{interviewbox}
